\documentclass{article}
\usepackage[UTF8]{ctex}
\usepackage{amsmath,amssymb,amsfonts}
\usepackage{enumitem}
\usepackage{geometry}
\geometry{a4paper, margin=1in}
% === 错题本排版优化（不参与题目标识） ===
\setlength{\parindent}{0pt}
\setlength{\parskip}{0.35em plus 0.10em minus 0.08em}
\setlength{\emergencystretch}{3em}
\setlist{itemsep=0.25em, topsep=0.40em, parsep=0pt}
\allowdisplaybreaks[2]
% === 错题本排版优化结束 ===
\begin{document}
\section*{湘潭大学2025年高等代数真题}
\section*{188．湘潭大学2025年高等代数试卷}
一．填空题．每题5分，共30分．\\

\subsection{题目1}
1．已知 \(A\) 为3阶方阵，且 \(|A-E|=|A-2024 E|=|A-2025 E|=0\) ，则 \(|A-2026 E|=\) \(\_\_\_\_\)．

\subsection{题目2}
2．已知多项式 \(f(x)=x^{3}+t x+2\) 有重根，则 \(t=\) \(\_\_\_\_\)．

\subsection{题目3}
3．设 \(V\) 是次数小于4的实系数多项式全体构成的线性空间，定义 \(\mathcal{D}(f(x))=f(x-1)\) ，则线性变换\(\mathcal{D}\) 在 \(1, x, x^{2}, x^{3}\) 这组基下的矩阵为 \(\_\_\_\_\)。

\subsection{题目4}
4．给定 \(A=\left(a_{i j}\right)\) 为2024阶方阵，\(a_{i j}=(i+j-1)^{2}\) ，则 \(|A|=\) \(\_\_\_\_\)．

\subsection{题目5}
5．已知 \(A\) 为2024阶可逆方阵，则 \(\left(\begin{array}{cc}O & A \\ A^{\prime} & O\end{array}\right)\) 的负惯性指数为 \(\_\_\_\_\)．\\
\(\_\_\_\_\)

\subsection{题目6}
6．与矩阵 \(\left(\begin{array}{ccc}1 & 2024 & 0 \\ 0 & 1 & 2024 \\ 0 & 0 & 1\end{array}\right)\) 乘法可交换的所有实矩阵在实数域上对矩阵加法和数乘构成一个线性空间，它的维数是 \(\_\_\_\_\)．

二．辨析题．正确的给出理由，错误的举出反例或说明理由．每题 6 分，共 36 分．\\

\subsection{题目7}
1．已知 \(B, C\) 为 \(n\) 阶方阵，\(E\) 为 \(n\) 阶单位方阵，则 \(|E-B C|=|E-C B|\) 。

\subsection{题目8}
2．若 \(n\) 阶方阵 \(A\) 的特征多项式有重根，则 \(A\) 不能相似对角化。

\subsection{题目9}
3．给定任意 \(n\) 阶实对称方阵 \(A\) ，存在实对称方阵 \(B\) ，使得 \(A=B^{2}\) 。

\subsection{题目10}
4．设 \(M_{1}, M_{2}, M_{3}\) 是线性空间 \(V\) 的子空间，若 \(M_{1} \cap M_{2}=M_{1} \cap M_{3}=M_{2} \cap M_{3}=\{0\}\) ，则子空间\(M_{1}+M_{2}+M_{3}\) 的和是直和。

\subsection{题目11}
5．设 \(V\) 是实数域上的一个2025维线性空间， \(\mathcal{A}\) 是 \(V\) 上任意的线性变换，则 \(V\) 有非平面的 \(\mathcal{A}\) 不变子空间。

\subsection{题目12}
6．存在 3 阶方阵 \(A\) ，使得 \(A\) 的伴随矩阵为 \(\left(\begin{array}{lll}1 & 1 & 1 \\ 1 & 1 & 1 \\ 1 & 1 & 1\end{array}\right)\) 。
\end{document}
